% Generated draft: requires mathematical and source review.
\documentclass[11pt]{article}
\usepackage[margin=25mm]{geometry}
\usepackage{fontspec}
\setmainfont{Latin Modern Roman}
\usepackage{amsmath,amssymb,mathtools,unicode-math}
\setmathfont{Latin Modern Math}
\usepackage{enumitem,needspace,xurl,hyperref}
\setlength{\parindent}{0pt}
\setlength{\parskip}{5pt}
\setlength{\emergencystretch}{3em}
\newcommand{\N}{\mathbb{N}}
\newcommand{\Z}{\mathbb{Z}}
\newcommand{\Q}{\mathbb{Q}}
\newcommand{\R}{\mathbb{R}}
\newcommand{\C}{\mathbb{C}}
\newcommand{\F}{\mathbb{F}}
\newcommand{\A}{\mathbb{A}}
\newcommand{\PP}{\mathbb{P}}
\newcommand{\E}{\mathbb{E}}
\newcommand{\eps}{\varepsilon}
\newcommand{\dd}{\,\mathrm{d}}
\begin{document}
\section*{UnsolvedMath ID: 30000486}
\textbf{Source problem code:} catalogue entry
\subsection{Definitions and mathematical statement}
An associative algebra is nil if every element is nilpotent; in this question the generated associative algebra must be understood without adjoining an identity. A quadratic algebra has a presentation by generators of degree one and homogeneous relations of degree two, equivalently a quotient of the positive-degree tensor algebra by an ideal generated in degree two. The first question asks whether such a nil algebra can have infinite dimension over its ground field. The source does not specify the ground field, its quantifier, or whether the space of generators is required to be finite dimensional.

For the two concrete questions let $F$ have characteristic $p>0$, and set
\[
R=F[t_0,t_1,\ldots]/(t_i^p:i\geq0),\qquad \partial_i=\frac{\partial}{\partial t_i}.
\]
The derivations descend to $R$ in characteristic $p$. With empty products equal to $1$, define
\[
V_1=\sum_{i\geq1}\left(\prod_{j=0}^{i-2}t_j^{p-1}\right)\partial_i,
\qquad
V_2=\sum_{i\geq2}\left(\prod_{j=1}^{i-2}t_j^{p-1}\right)\partial_i.
\]
These sums act locally finitely on every polynomial of $R$. Let $A\subseteq\operatorname{End}_F(R)$ be the $F$-linear span of all nonempty words in $V_1,V_2$, with multiplication given by composition. The second question is
\[
\forall a\in A\quad\exists n\geq1:\quad a^n=0\;?
\]
Let $L\subseteq\operatorname{Der}_F(R)$ be the Lie algebra generated by $V_1,V_2$, with commutator bracket, let $U=U(L)$ be its ordinary universal enveloping algebra, and let $Z=Z(U)$. The third question asks whether
\[
D=(Z\setminus\{0\})^{-1}U
\]
is an algebraic division algebra. Here the assertion that $D$ is a division algebra means that every nonzero element is invertible. A natural possible meaning of algebraicity is that every $d\in D$ satisfies a nonzero polynomial over $K=\operatorname{Frac}(Z)$; algebraicity over $F$ is a different requirement. The supplied formulation does not identify the base field for algebraicity, so this distinction remains unresolved. The first existence question and the two questions about this construction are retained separately; no quadratic presentation of $A$ is supplied.
\subsection{Short English statement}
Ask whether an infinite-dimensional quadratic nil algebra exists, whether the specified two derivations generate a nil associative algebra, and whether the central localization of their Lie algebra's enveloping algebra is an algebraic division algebra. The algebraicity base field and some conventions for the first question require clarification.
\subsection{Sources}
\begin{itemize}
\item Supplied report: \emph{Interactions between Algebraic Geometry and Noncommutative Algebra} (2007), Oberwolfach Reports, DOI 10.4171/owr/2006/23, pp. 1319--1384; proposed year 2006. \url{https://doi.org/10.4171/owr/2006/23}.
\item Supplied literature reference, labelled ``recent Kurosh-problem survey and results'': \url{https://www.cambridge.org/core/journals/forum-of-mathematics-sigma/article/quantitative-kurosh-problem/F66BF901CF6A811926EAC2D30C17B409}.
\item Local UnsolvedMath registry, permanent ID 30000486, source code catalogue entry; created 25 August 2026, updated 16 September 2026. The record labels the statement corrected\_verified and reports comparison with the official report. Its literature triage of 21 August 2026 labels the three-question record open and reports no general resolution of the quadratic nil-algebra or division-algebra questions in the checked literature; this is a supplied dated assessment requiring independent verification.
\end{itemize}
\end{document}
